\documentclass[10pt]{letter}
\usepackage[a4paper, margin=0.85in]{geometry}
\usepackage[T1]{fontenc}
\setlength{\parskip}{5pt}
\signature{Harsh Bordekar \\ harshbordekar@gmail.com}
\begin{document}
\begin{letter}{Airbus Operations GmbH \\ Hamburg}
\opening{Dear Hiring Team,}

I am writing to apply for the F\&DT Structural Analysis Engineer position in Hamburg. I am a structural analysis and simulation engineer with over seven years in aerospace structural integrity, and building and validating structural analysis models against defined methods, then standing behind them under technical review, is the core of what I already do.

At DLR I regularly verify and technically review externally and student-produced analysis models, meshes, and results against defined methods and acceptance criteria, providing sign-off-style checks on structural analysis data and documentation and communicating findings to design and project teams. During my PhD I derived a certification-aligned failure criterion that translates multiscale FEM outputs into structural sizing inputs consistent with regulatory boundary conditions, the same kind of regulatory-boundary thinking your posting frames as applying airworthiness regulations. I also build the automation that this posting calls digital transformation: Python- and HPC-based automation of analysis and reporting, and an agentic AI application that predicts material properties on demand instead of requiring a full-scale FEM run each time.

I work daily with the FE solvers (Nastran/Patran, ABAQUS) and analytical stress methods (Niu, Bruhn) your posting names, and with progressive damage analysis of composites (CDM, PDA, CZM) alongside metallic fatigue and damage tolerance work using small-defect fatigue models (Murakami, El-Haddad). I hold an M.Sc. in Computational Methods in Engineering and I am comfortable working autonomously and communicating with stakeholders across design and project teams.

To be transparent about fit: my structural analysis work has been on hydrogen storage tanks, a rocket housing, and additively manufactured titanium parts, not on the A350XWB or A320 Family specifically, and I have not used ISAMI or HSB, only Niu and Bruhn. I have also not personally defined physical structural test campaigns for certification; my closest real experience is coupling defect measurements with numerical simulation to set qualification and acceptance criteria. I would rather name these directly than have them surface later, and I am confident the underlying structural analysis, review, and regulatory-compliance discipline transfer quickly to Airbus's aircraft programs.

I am motivated by the chance to bring rigorous, validated structural analysis and a genuine appetite for automating engineering workflows to Airbus, and I would welcome the opportunity to discuss how I can contribute to your team.

Thank you for your consideration. I look forward to discussing how I can contribute.

\closing{Sincerely,}
\end{letter}
\end{document}
